\documentclass[11pt,letterpaper]{article}
\usepackage[margin=0.88in]{geometry}
\usepackage[T1]{fontenc}
\usepackage{lmodern,microtype,amsmath,amssymb,bm,booktabs,longtable,array,graphicx}
\usepackage[hidelinks]{hyperref}
\usepackage{enumitem}
\setlist{nosep,leftmargin=1.4em}
\setlength{\parindent}{0pt}
\setlength{\parskip}{5pt}
\setlength{\emergencystretch}{2em}
\newcommand{\J}{\mathcal J}
\newcommand{\dd}{\mathrm d}
\newcommand{\X}{\mathbf X}
\newcommand{\vct}[1]{\boldsymbol{#1}}
\newcommand{\R}{\mathbb R}
\newcommand{\tr}{\operatorname{tr}}
\newcommand{\diag}{\operatorname{diag}}
\newcommand{\norm}[1]{\left\lVert #1\right\rVert}
\title{Magnetic-moment-constrained collisions\\\large Equations, numerical design and verification notes}
\author{Project: \texttt{sato-morrison}}
\date{6 October 2026}
\begin{document}
\maketitle
\begin{abstract}
These notes specify a prototype based on Sato and Morrison's reduced collision
formalism. They distinguish the published equations, derived test cases, proposed
numerical methods and unresolved physical assumptions. The supplied executable
code verifies a uniform-field linear limit only. Nonlinear evolution, nonuniform
geometry and dimensional collision-rate validation remain implementation work.
\end{abstract}

\section{Scope and equation map}

The primary source is Ref.~\cite{SM25}. Its Eq.~(35) gives a one-species
noncanonical collision structure. Equations (111) and (119) are different levels
of description: the former evolves a ten-coordinate pair distribution, while the
latter closes two five-coordinate species distributions after a pair-factorization
assumption. Equation (181) is a localized linear approximation with a prescribed
scalar collision coefficient. This prototype must identify which of these models
it implements.

The starting model has a prescribed, nonzero vacuum magnetic field, no magnetic
perturbation, a fixed electrostatic reference potential, and one species. A
nonlinear local-kernel extension will be named separately from literal Eq.~(181).
No absolute collision rate or laboratory confinement result is established by the
starter calculation.

\begin{center}
\begin{tabular}{p{0.22\linewidth}p{0.67\linewidth}}
\toprule
Source equations & Implementation role\\
\midrule
SM25 (35), (39)--(59) & Weak form, boundary conditions, invariants and stationary family.\\
SM25 (73)--(100) & Guiding-center bracket, vacuum assumption and invariant measure.\\
SM25 (118)--(119) & Pair closure and multispecies kernel; later stage.\\
SM25 (139)--(150) & Constrained stationary states and initial magnetic-moment marginal.\\
SM25 (164)--(183) & Linearization, kernel approximations and the first numerical target.\\
BS25, Sec.~V & Structural comparison to a guiding-center Landau bracket.\\
\bottomrule
\end{tabular}
\end{center}

\section{Coordinates, measures and the bracket}

We use SI variables before nondimensionalization. Let $u$ be parallel velocity,
$\mu=m v_\perp^2/(2B)$ the magnetic moment, and $\eta=\mu B$ the perpendicular
kinetic energy. With $\bm b=\bm B/B$,
\begin{equation}
 \bm B^*=\bm B+\frac{m u}{q}\nabla\times\bm b,
 \qquad B_\parallel^*=\bm b\cdot\bm B^*.
 \label{eq:bstar}
\end{equation}
For the vacuum backgrounds used in Ref.~\cite{SM25},
$\bm b\cdot\nabla\times\bm b=0$, hence $B_\parallel^*=B$. This does not imply
$\nabla\times\bm b=0$: curvature terms in $\bm B^*$ generally remain.

In coordinates $z=(\X,u,\mu)$ the gyrophase-independent bracket acts on a
covector $a=(\bm a_X,a_u,a_\mu)$ as
\begin{equation}
 \J_\mu a=
 \begin{pmatrix}
 (\bm b\times\bm a_X)/(qB)+\bm B^*a_u/(mB)\\
 -\bm B^*\cdot\bm a_X/(mB)\\0
 \end{pmatrix}.
 \label{eq:jmu}
\end{equation}
The Hamiltonian and characteristic equations are
\begin{align}
 E&=\frac12m u^2+\mu B+q\Phi,\label{eq:energy}\\
 \dot\X&=u\frac{\bm B^*}{B}
 +\frac{\bm b\times(\mu\nabla B+q\nabla\Phi)}{qB},\label{eq:xdot}\\
 \dot u&=-\frac{\bm B^*\cdot(\mu\nabla B+q\nabla\Phi)}{mB},
 \qquad \dot\mu=0.\label{eq:udot}
\end{align}
The signs in \eqref{eq:jmu} must be checked against the cross-field drift in
\eqref{eq:xdot}, not inferred from a memorized matrix convention.

The phase-space measures are
\begin{equation}
 \dd\Gamma=B\,\dd^3X\,\dd u\,\dd\mu
 =\dd^3X\,\dd u\,\dd\eta.
 \label{eq:measure}
\end{equation}
We follow the article's normalization: $n(\X)=\int f\,\dd u\,\dd\eta$. If $F$
is an ordinary gyrotropic velocity density normalized with $\dd^3v$, then
$\dd^3v=(2\pi B/m)\dd u\dd\mu$ and $f=(2\pi/m)F$. Track this constant,
particularly in multispecies comparisons.

The coordinate change $\zeta=(\X,u,\eta)$ has derivative
\begin{equation}
 T=\frac{\partial\zeta}{\partial z}
 =\begin{pmatrix}I_3&0&0\\0&1&0\\
 \mu(\nabla B)^T&0&B\end{pmatrix},
 \qquad \J_\eta=T\J_\mu T^T.
 \label{eq:transform}
\end{equation}
This constructs the matrix in SM25 (93), equivalently (182), without duplicating
its entries by hand. It also gives
\begin{equation}
 \J_\eta\nabla_\zeta(\eta/B)=0,
 \qquad \dot\eta=\mu\dot\X\cdot\nabla B.
 \label{eq:etadot}
\end{equation}
Thus a discretization that freezes $\eta$ during spatial transport is not the
magnetic-moment-preserving model in a nonuniform field.

With general coordinates and measure $\rho(z)\dd z$, conservative divergence
means $\rho^{-1}\partial_a(\rho F^a)$, not $\partial_aF^a$. Verify the Liouville
identity $\partial_a(\rho\J^{ab})=0$. The prototype uses $\rho=B$ in the $\mu$
chart and $\rho=1$ in the article's $\eta$ chart.

\subsection*{Normalization is part of the closure}

Choose fixed scales $L_*,v_*,E_*,t_*$ and a diagonal coordinate scale matrix $S$.
For $\zeta=S\hat\zeta$, $E=E_*\hat E$ and $t=t_*\hat t$,
\begin{equation}
 \hat\J=t_*E_*S^{-1}\J S^{-T}.
 \label{eq:dimensionless}
\end{equation}
State these scales with every dimensional result. A Euclidean phase-space
projector applied to a vector containing spatial velocities, accelerations and
energy rates is not defined independently of those scales.

For a purely affine chart change $y=A z$ with time and energy conventions fixed,
a pair vector transforms as $d_y=A d_z$. Its quadratic form is unchanged only if
\begin{equation}
 \Pi_y=A^{-T}\Pi_zA^{-1},\qquad
 d_y^T\Pi_y e_y=d_z^T\Pi_z e_z.
 \label{eq:kerneltransform}
\end{equation}
Recreating an identity matrix and Euclidean projector in the new chart generally
specifies a different kernel. For the position-dependent map \eqref{eq:transform},
compute the two vectors in one common chart before subtracting:
\begin{equation}
 \Delta_\eta\phi
 =T(z)\J_\mu(z)\nabla_z\phi
 -T(z')\J_\mu(z')\nabla_{z'}\phi'.
 \label{eq:commonchart}
\end{equation}
This is the safe route when storing the state on $\mu$ nodes but reproducing a
kernel defined in the $\eta$ chart.

\section{Nonlinear collision structure and conservation}

For a single species define
\begin{equation}
 \Delta\phi=\J\nabla\phi-\J'\nabla'\phi',\qquad \xi=\Delta E.
 \label{eq:delta}
\end{equation}
A specified kernel must obey exchange symmetry, symmetry of its quadratic form,
positive semidefiniteness and the energy degeneracy:
\begin{equation}
 \Pi=\Pi^T\succeq0,\qquad \Pi\xi=0.
 \label{eq:kernel}
\end{equation}
The source strong form is
\begin{equation}
 C[f]=\nabla_z\cdot\left[
 f\J\int f'\Pi(\J'\nabla'\ln f'-\J\nabla\ln f)\,\dd z'\right]
 \label{eq:strong}
\end{equation}
in a flat invariant-measure chart. Integration by parts, antisymmetry of $\J$
and exchange of the two labels give
\begin{equation}
 \int\phi C[f]\dd\Gamma
 =-\frac12\iint ff'\,\Delta\phi^T\Pi\Delta\ln f\,
 \dd\Gamma\dd\Gamma'.
 \label{eq:weak}
\end{equation}
All following equalities require the boundary terms to vanish. The paper states
conditions for both the total flux and the Hamiltonian part; a finite numerical
box does not satisfy them automatically.

Set $\phi=1$ in \eqref{eq:weak} to obtain particle conservation. Set $\phi=E$
and use \eqref{eq:kernel} for collision energy conservation. Since
$\J\nabla g(\mu)=0$ for every smooth $g$,
\begin{equation}
 \frac{\dd}{\dd t}\int f g(\mu)\dd\Gamma=0.
 \label{eq:casimirs}
\end{equation}
For closed evolution this is equivalent to conserving the full marginal
\begin{equation}
 G(\mu)=\int B f(\X,u,\mu)\dd^3X\dd u.
 \label{eq:marginal}
\end{equation}
It is stronger than conservation of $\int\mu f\dd\Gamma$ alone.

Define $S[f]=-\int f\ln(f/f_{\rm ref})\dd\Gamma$ with a constant reference
scale inside the logarithm. Particle conservation removes that scale from the
entropy rate. Equation \eqref{eq:weak} yields
\begin{equation}
 \dot S=\frac12\iint ff'\Delta\ln f^T\Pi\Delta\ln f\,
 \dd\Gamma\dd\Gamma'\ge0.
 \label{eq:entropy}
\end{equation}
These are structural statements from the model, not a proof of a unique final
state. Additional null directions of $\Pi$ can prevent mixing.

A candidate observable $p$ is conserved by the full kinetic model only if the
Hamiltonian flow preserves it and the collision form annihilates $\Delta p$.
An external dipole field does not imply conservation of all Cartesian momenta.
For an axisymmetric model, check canonical angular momentum using the actual
kernel and boundaries rather than inferring it from the field name.

\section{The fixed-field linear model}

Write $f=f_0(1+h)$, where $f_0>0$ is stationary. For the source stationary family,
\begin{equation}
 f_0=A\exp[-\beta E_0+g(\mu)],
 \qquad \J\nabla\ln f_0=-\beta\J\nabla E_0.
 \label{eq:fzero}
\end{equation}
With the fields and kernel fixed, the zeroth-order collision term vanishes.
The first-order weak form is
\begin{equation}
 \int\phi C_L[f_0h]\dd\Gamma
 =-\frac12\iint f_0f_0'\Delta\phi^T\Pi_0\Delta h\,
 \dd\Gamma\dd\Gamma'.
 \label{eq:linearweak}
\end{equation}
The local scalar approximation in SM25 (181) uses
\begin{equation}
 \Pi_0=D P I_x P\,\delta^3(\X-\X'),\quad
 P=I-\frac{\xi_0\xi_0^T}{\xi_0^T\xi_0},\quad
 I_x=\diag(1,1,1,0,0).
 \label{eq:local}
\end{equation}
Here $D$ is a prescribed constant, and the projector is evaluated in a specified
normalized common chart. The two velocity integrals remain after the spatial
delta function is used. In a $\mu$ chart, the local pair measure contains two
velocity Jacobians; do not accidentally retain only one factor of $B$.

The perturbation norm
\begin{equation}
 Q[h]=\frac12\int f_0 h^2\dd\Gamma
 \label{eq:q}
\end{equation}
satisfies
\begin{equation}
 \dot Q=-\frac12\iint f_0f_0'\Delta h^T\Pi_0\Delta h\,
 \dd\Gamma\dd\Gamma'\le0
 \label{eq:qrate}
\end{equation}
for collision-only dynamics. A compatible stationary Hamiltonian flow also
preserves $Q$ under suitable boundary conditions.

For self-consistent fields, \eqref{eq:linearweak} is not a complete linearization
unless the omitted terms have been ordered out. Differentiating the exact
energy-degeneracy condition gives
\begin{equation}
 \delta\Pi\,\xi_0+\Pi_0\delta\xi=0.
 \label{eq:deltapi}
\end{equation}
The approximations in SM25 (164)--(170) must be carried explicitly into any
extension. Compare the proposed linear operator with a directional derivative of
the nonlinear residual at fixed physical normalization.

\section{Uniform-field reduction and exact null modes}

This section is an independent derivation from the source Eq.~(181). Take a
uniform field, a constant reference potential and $f_0=f_0(u,\eta)$. For
$\bm b=\bm e_z$,
\begin{equation}
 \J a=\left(-\frac{a_y}{qB},\frac{a_x}{qB},
 \frac{a_u}{m},-\frac{a_z}{m},0\right)^T,
 \qquad \J\nabla E_0=(0,0,u,0,0)^T.
 \label{eq:uniformj}
\end{equation}
Consequently $\xi_0=(u-u')\bm e_z$ in the spatial block, and
\begin{equation}
 P I_xP=\diag(1,1,0,0,0).
 \label{eq:uniformp}
\end{equation}
At $u=u'$ use the continuous limit \eqref{eq:uniformp}; this prescription does not
extend automatically to a general nonuniform field.

Only perpendicular spatial derivatives survive. Let
\begin{equation}
 n_0=\int f_0\dd u\dd\eta,\qquad
 \delta n=\int\delta f\dd u\dd\eta.
\end{equation}
Two applications of the cross-product matrix give
$(\bm b\times)^2=-I_\perp$. Substitution into the strong form gives
\begin{equation}
 \boxed{C_L[\delta f]=\frac{D}{(qB)^2}\nabla_\perp^2
 (n_0\delta f-f_0\delta n).}
 \label{eq:uniformreduction}
\end{equation}
This is spatial mixing of velocity-dependent structure, not ordinary homogeneous
velocity-space thermalization.

For a Fourier mode with $\bm k_\perp=\bm k-(\bm k\cdot\bm b)\bm b$, define
\begin{equation}
 \lambda=\frac{Dn_0|\bm k_\perp|^2}{(qB)^2},\qquad
 \bar h=\frac{1}{n_0}\int f_0h\dd u\dd\eta.
\end{equation}
The collision-only equation is
\begin{equation}
 \partial_t h=-\lambda(h-\bar h),\qquad
 h(t)=\bar h+[h(0)-\bar h]e^{-\lambda t}.
 \label{eq:exactdecay}
\end{equation}
Thus a density-like perturbation is undamped. If $\bm k_\perp=0$, every velocity
perturbation is collisionally undamped. A homogeneous anisotropic distribution
is therefore not an isotropization benchmark for this model.

With positive velocity quadrature weights $w_i$, set
\begin{equation}
 m_i=w_i f_{0i},\quad s_i=\sqrt{m_i},\quad n_0=\sum_i m_i.
\end{equation}
The matrix for $h$ and its entropy-scaled form are
\begin{align}
 L_h&=-\lambda\left(I-\frac{\bm 1\bm m^T}{n_0}\right),\\
 M_v^{1/2}L_hM_v^{-1/2}
 &=-\lambda\left(I-\frac{\bm s\bm s^T}{n_0}\right).
 \label{eq:uniformspectrum}
\end{align}
For $\lambda>0$ the spectrum has one zero eigenvalue and $N_v-1$ copies of
$-\lambda$. This is an exact finite-quadrature statement when the same $n_0$ is
used throughout. It provides a stringent factor, sign and normalization test.

When parallel streaming is added, $\partial_t h=-ik_\parallel uh+L_hh$.
A density-like collision null mode need not remain in that nullspace under
streaming. Distinguish the collision spectrum from the spectrum of the combined
generator. With $\bm k_\perp=0$ and frozen fields the exact streaming solution is
$h(t,u,\eta)=h(0,u,\eta)e^{-ik_\parallel ut}$; this is not self-consistent Landau
damping.

\section{Stationary states and initial constraints}

The source stationary family \eqref{eq:fzero} follows from energy degeneracy and
the magnetic-moment Casimir. For $g(\mu)=-\beta\gamma\mu$,
\begin{equation}
 f_*=A\exp\left[-\beta\left(\frac12m u^2+
 \mu(B+\gamma)+q\Phi\right)\right],\quad B+\gamma>0.
 \label{eq:linearg}
\end{equation}
Integrating over $u$ and $\mu$ using \eqref{eq:measure} gives
\begin{equation}
 n_*(\X)=A\sqrt{\frac{2\pi}{\beta m}}\,
 e^{-\beta q\Phi}\frac{B}{\beta(B+\gamma)}.
 \label{eq:density}
\end{equation}
The corresponding temperature parameters in energy units are
\begin{equation}
 T_\parallel=\beta^{-1},\qquad
 T_\perp=\frac{B}{\beta(B+\gamma)}.
 \label{eq:temperatures}
\end{equation}
These formulas are useful checks of velocity integration and the definition of
pressure. They do not establish that an arbitrary initial state evolves to this
particular family member.

For a prescribed fixed-field energy and a known initial marginal $G_0(\mu)$,
let
\begin{equation}
 Z_\mu(\beta)=\int B(\X)e^{-\beta E(\X,u,\mu)}\dd^3X\dd u.
\end{equation}
A candidate constrained maximum-entropy state is
\begin{equation}
 f_*(\X,u,\mu)=\frac{G_0(\mu)}{Z_\mu(\beta)}e^{-\beta E(\X,u,\mu)}.
 \label{eq:generalsteady}
\end{equation}
Solve for $\beta$ from total energy. Add other invariant multipliers if they
actually survive the chosen dynamics. The same initial inverse-temperature
parameter need not equal the final multiplier. Extra kernel invariants can make
\eqref{eq:generalsteady} inaccessible.

For self-consistent electrostatics,
\begin{equation}
 -\epsilon_0\nabla^2\Phi=\sum_s q_sn_s+\rho_{\rm fixed}.
 \label{eq:poisson}
\end{equation}
The self-field energy is $\epsilon_0\int|\nabla\Phi|^2\dd^3X/2$, with boundary
terms treated consistently. A prescribed external potential contributes
$\sum_s\int q_s\Phi_{\rm ext}f_s\dd\Gamma_s$ without the self-interaction
half-factor. A periodic field solve requires net neutrality and a gauge choice.

\section{Comparison models and what each tests}

The conventional reduced Landau literature retains magnetic-moment friction and
diffusion even when the distribution is gyrophase independent
\cite{B04,BBQ15,HBP17,BS25}. Brizard and Sugama explicitly discuss a structural
reduction with $\bm\Delta_{\rm gc}^{\mu}=0$ and a reduced relative velocity in
Sec.~V of Ref.~\cite{BS25}. This is a bridge for comparison, not a verification
that every simplified Sato--Morrison kernel is identical to their Landau tensor.
The papers use different conventions; convert all masses, charges and unit
factors before comparing implementations.

\subsection*{Pitch-angle scattering}

With $v=|\bm v|$ and $\xi=u/v$, use
\begin{equation}
 C_{\rm Lz}[F]=\frac{\nu(v)}{2}\partial_\xi
 [(1-\xi^2)\partial_\xi F].
 \label{eq:lorentz}
\end{equation}
The Legendre identity gives $C_{\rm Lz}[P_\ell]=
-\nu\ell(\ell+1)P_\ell/2$ for constant $\nu$ on a speed shell. Hence degrees
$0,1,2,3$ have rates $0,\nu,3\nu,6\nu$. Speed-shell population and energy are
conserved, but momentum of the evolved species can be transferred to the fixed
reservoir. Since $\mu=m v^2(1-\xi^2)/(2B)$, this is a moment-changing comparison.
Derive its coordinate transform; it is not bare diffusion at fixed $u$ in $\mu$.

\subsection*{A nonlinear conserving Fokker--Planck model}

The single-species Dougherty form used as a comparison is
\begin{equation}
 C_D[F]=\nu\nabla_v\cdot[(\bm v-\bm U)F+\theta\nabla_vF],
 \quad \bm U=\frac1n\int\bm vF\dd^3v,
 \quad \theta=\frac{1}{3n}\int|\bm v-\bm U|^2F\dd^3v.
 \label{eq:dougherty}
\end{equation}
Integration by parts shows number and momentum conservation. The trace of the
centered covariance is constant. Its matrix satisfies
\begin{equation}
 \dot\Theta=-2\nu(\Theta-\theta I),\qquad
 \Theta(t)=\theta I+[\Theta(0)-\theta I]e^{-2\nu t}.
 \label{eq:covariance}
\end{equation}
For an initially Gaussian homogeneous state this specifies the full Gaussian
solution. Use it as an exact code test. A fixed bath temperature or drift would
define a different model. Neither version is the exact Coulomb operator.

\subsection*{Physical three-velocity Landau reference}

A one-species normalized reference is
\begin{align}
 U(\bm w)&=\frac{I-\hat{\bm w}\hat{\bm w}^{T}}{|\bm w|},\\
 C_{\rm La}[F]&=\Gamma\nabla_v\cdot\int FF'U(\bm v-\bm v')
 (\nabla_v\ln F-\nabla_{v'}\ln F')\dd^3v'.
 \label{eq:landau}
\end{align}
Derive the dimensional prefactor for the selected normalization separately. Its
positive kernel and velocity-difference null direction support conventional
number, momentum, energy and entropy tests. Use an independently assembled
small 3V problem before implementing a general transformed operator.

A gyrotropic two-coordinate distribution still has three-dimensional collision
geometry. Converge the gyrophase integral or derive its analytic average; a
literal 2V Landau kernel is not an equivalent comparator. Maxwell-molecule and
BKW-type tests in Ref.~\cite{CHWW20} can test a different kernel's numerical
machinery, but do not establish Coulomb rates.

\section{Analytic magnetic fields and boundaries}

A vacuum field can be specified locally as the gradient of a harmonic scalar
potential. Use the following fields with explicit domains.

\textbf{Uniform field.} $\bm B=B_c\bm b_0$, with $|\bm b_0|=1$.

\textbf{Harmonic mirror.} Let $R^2=x^2+y^2$ and
\begin{align}
 \Psi_B&=B_cz+a\left(z^3/3-zR^2/2\right),\\
 \bm B&=\nabla\Psi_B=(-azx,-azy,B_c+a[z^2-R^2/2]).
 \label{eq:mirrorfield}
\end{align}
Its curl and divergence vanish identically. Choose a domain on which $B$ stays
strictly positive. The dimensionless variation is controlled by $aL^2/B_c$.
An arbitrary field $B(z)\bm e_z$ is not an admissible substitute.

\textbf{Toroidal vacuum field.} On $R>R_{\min}>0$,
\begin{equation}
 \bm B=\frac{C}{R^2}(-y,x,0).
 \label{eq:toroid}
\end{equation}
This provides curvature and field-strength variation without a parallel gradient
of $B$. Its noncontractible domain is part of the problem.

\textbf{Point dipole.} Away from the source,
\begin{equation}
 \bm B=\frac{C_d}{r^5}(3xz,3yz,2z^2-x^2-y^2),
 \quad \Psi_B=-C_dz/r^3.
 \label{eq:dipole}
\end{equation}
This is a field test, not a model of a finite winding or material boundary. Exclude
$r=0$ and any near-source region outside the guiding-center ordering.

\textbf{Symmetry breaking.} A controlled addition
$\delta\Psi_B=\varepsilon B_cxy/L$ is harmonic. Its gradient can be added to a
vacuum background; check the resulting nonzero-field domain and boundary setup.

Do not make a nonperiodic magnetic field periodic by wrapping its array. Begin
with collision-only nonuniform tests under natural zero-collision-flux boundary
conditions. For full dynamics, require tangent Hamiltonian flux, a derived
invariant-compatible reflection, a specified confining potential, or an open
system with measured flux budgets. Arbitrary specular reflection need not
preserve $\mu$. A density-profile test without these choices is not a closed
relaxation experiment.

\section{Linear weak discretization and SOLVAX}

Represent $h=\sum_i a_i\chi_i$. Define the mass and collision matrices by
\begin{align}
 M_{ij}&=\int f_0\chi_i\chi_j\dd\Gamma,\\
 K_{ij}&=\frac12\iint f_0f_0'\Delta\chi_i^T\Pi_0\Delta\chi_j
 \dd\Gamma\dd\Gamma'.
 \label{eq:gram}
\end{align}
Positive quadrature gives a Gram form $K=B_p^TW_pB_p\succeq0$. Implement its
action without assembling the full pair tensor. A dense assembly on a tiny grid
must remain an independent reference.

If $c$ represents a collision invariant and the derivative/kernel construction
respects it, then $Kc=0$. Include the discrete energy and all resolved
magnetic-moment-bin constraints. Finite differences do not automatically obey
the product rules needed for an arbitrary sampled $\mu B$. Use the same discrete
energy gradient in the degeneracy direction and in the weak operator, and prove
consistency as the mesh is refined.

The collision-only system and norm balance are
\begin{equation}
 M\dot a=-Ka,\qquad
 \frac{\dd}{\dd t}\frac12a^TMa=-a^TKa\le0.
 \label{eq:semidiscrete}
\end{equation}
With a compatible stationary Hamiltonian discretization,
\begin{equation}
 M\dot a=(A-K)a,\qquad A^T=-A.
 \label{eq:combined}
\end{equation}
Also test $Ac=0$ for the required linear invariants; skew symmetry alone does not
imply those individual identities.

A backward-Euler collision step solves
\begin{equation}
 (M+\Delta tK)a^{n+1}=Ma^n.
 \label{eq:backward}
\end{equation}
For positive-definite $M$, this is an appropriate PCG problem. SOLVAX's
\texttt{pcg\_linear\_solve} supplies a converged-system differentiation wrapper
\cite{SOLVAX}. Do not pass $M^{-1}K$ to Euclidean PCG without mass scaling or a
proof of symmetry. A coupled implicit step with $A$ is generally nonsymmetric
and calls for GMRES or another suitable method.

For a linear invariant $c$, multiply \eqref{eq:backward} by $c^T$ to check its
conservation. Backward Euler also satisfies
\begin{equation}
 Q^{n+1}-Q^n=-\Delta t(a^{n+1})^TKa^{n+1}
 -\frac12\norm{a^{n+1}-a^n}_M^2\le0.
\end{equation}
This does not prove positivity of a reconstructed finite-amplitude distribution.
An implicit midpoint step can be second order, but is not automatically a good
stiff-decay scheme or a nonlinear entropy-preserving integrator.

Local pair application costs $O(N_XN_v^2)$ in a direct quadrature. A spatially
nonlocal direct application can cost $O(N_X^2N_v^2)$. Chunk pair contractions,
reuse geometry and quadrature, and keep derivatives sparse. Block-tridiagonal
solves may be useful in a line preconditioner; the full velocity coupling is not
generally tridiagonal.

\section{A nonlinear discrete entropy construction}

This is a proposed finite-dimensional construction, not an implemented guarantee
for a general spatial mesh. It follows the weak-structure approach developed for
Landau discretizations \cite{KH17,HBK18,H21,ZPH}.

Let $f_i>0$, $w_i>0$, and $n_i=w_if_i$. Construct a symmetric positive mobility
$K(f)$ from the discrete pair weak form, with
\begin{equation}
 K\bm1=0,\quad Ke=0,\quad Kc_\mu=0,
 \label{eq:discreteconstraints}
\end{equation}
where $e$ is the fixed-field energy vector and $c_\mu$ spans the represented
moment-marginal constraints. Define
\begin{equation}
 S(n)=-\sum_i n_i\ln(n_i/w_i).
\end{equation}
Then the collision ODE can be written
\begin{equation}
 \dot n=K(f)\nabla_nS=-K(f)\ln f.
 \label{eq:nonlinearode}
\end{equation}
The constant term in $\nabla_nS=-\ln f-\bm1$ vanishes through
\eqref{eq:discreteconstraints}. Consequently $\dot S\ge0$ and
$\dd(e^Tn)/\dd t=0$. This is an algebraic guarantee only for the actual
constructed $K$.

For each component define a divided difference of $s_i(n)=-n\ln(n/w_i)$:
\begin{equation}
 (\overline\nabla S)_i=
 \begin{cases}
 [s_i(n_i^+)-s_i(n_i^-)]/(n_i^+-n_i^-),&n_i^+\ne n_i^-,\\
 -\ln(n_i^-/w_i)-1,&n_i^+=n_i^-.
 \end{cases}
 \label{eq:discretegradient}
\end{equation}
Evaluate it stably near equal arguments using logarithmic divided differences,
not cancellation-prone subtraction. A step
\begin{equation}
 n^+-n^-=\Delta t\,\overline K\,\overline\nabla S,
 \qquad \overline K=\overline K^T\succeq0
 \label{eq:discretestep}
\end{equation}
with the null directions in \eqref{eq:discreteconstraints} gives
\begin{equation}
 S(n^+)-S(n^-)=\Delta t\,
 (\overline\nabla S)^T\overline K\overline\nabla S\ge0.
\end{equation}
The conservation laws follow by left multiplication with each invariant vector.
Use a positive mean state to construct $\overline K$ and a converged nonlinear
solve. Parameterizing $n_i^+=w_i e^{g_i^+}$ enforces positive nodal populations,
but does not prove existence of the solve or positivity between interpolation
nodes. Check all required quadrature/reconstruction locations.

For self-consistent fields the energy is nonlinear in $n$. Its discrete gradient
must enter the collision degeneracy consistently; the fixed vector $e$ above is
not enough. Likewise, nonlinear Hamiltonian transport requires its own
conservative and positivity-compatible treatment. A linear skew matrix alone
is not a proof of a nonlinear H-theorem.

\section{Scattering physics and validity tests}

The constrained model assumes negligible accumulated change of magnetic moment.
That assumption is stronger than a small collision frequency compared with the
gyrofrequency. At constant $B$, an energy change satisfies
\begin{equation}
 \Delta E=\frac{m}{2}\Delta(u^2)+B\Delta\mu+q\Delta\Phi.
 \label{eq:energychange}
\end{equation}
Small $\Delta E$ does not separately bound $\Delta\mu$. For independent unbiased
increments, zero mean accumulation can coexist with
\begin{equation}
 \operatorname{Var}\left(\sum_{j=1}^N\delta\mu_j\right)
 =N\operatorname{Var}(\delta\mu).
\end{equation}
These are audit points for the physical ordering around SM25 (85)--(86), not
replacements for that paper's stated assumptions.

A controlled microscopic reference can integrate
\begin{equation}
 m_s\dot{\bm v}_s=q_s\bm v_s\times\bm B-\nabla_{\bm r_s}V,
 \qquad \dot{\bm r}_s=\bm v_s,
 \label{eq:encounter}
\end{equation}
with a specified screened potential, for example
$V(r)=q_1q_2e^{-r/\lambda}/(4\pi\epsilon_0r)$. Screening is a declared model
choice, not a result of this two-body calculation. Resolve close encounters or
exclude them with a stated grazing criterion. The generalized magnetized
scattering theory in Ref.~\cite{JB20} is an important reference.

For uniform $\bm B$, compare with the unperturbed helical orbit at identical
incoming conditions. Use
\begin{equation}
 \bm R_s=\bm r_s+\frac{\bm v_s\times\bm b}{\Omega_s},\qquad
 \mu_s=\frac{m_s|\bm v_{\perp s}|^2}{2B},\quad \Omega_s=q_sB/m_s.
\end{equation}
Extract asymptotic incoming/outgoing guiding-center data at a common reference
time. A raw endpoint displacement can contain an arbitrary flight-time
contribution after the encounter changes parallel velocity.

Sample incident flux, impact parameters, gyrophases and relative velocities with
the appropriate measure. Check timestep, start/end distance, angular sampling and
screening/cutoff convergence independently. Recover a Rutherford limit before
using strong-field results. Attractive close encounters outside the grazing
regime are not automatically described by the reduced kinetic closure.

For independent encounters at a stated rate $\nu_{\rm enc}$, diagnostic drift and
diffusion coefficients are
\begin{equation}
 A_\mu=\nu_{\rm enc}\langle\delta\mu\rangle,\qquad
 D_{\mu\mu}=\frac{\nu_{\rm enc}}2\langle(\delta\mu)^2\rangle.
 \label{eq:driftdiffusion}
\end{equation}
Use both, together with cross-covariances and correlation tests. A timescale
estimate must specify a nonzero population scale $\mu_{\rm ref}$; the quantities
$\mu_{\rm ref}/|A_\mu|$ and $\mu_{\rm ref}^2/(2D_{\mu\mu})$ diagnose different
failures of the constraint. Near $\mu=0$, use a suitable distribution-level metric.

Compare these with an independently determined mixing time. A physically
supported constrained interval requires
\begin{equation}
 \tau_{\rm mix}\ll t\ll\tau_\mu.
\end{equation}
Choosing a small input pitch-angle rate establishes a synthetic asymptotic test,
not this physical inequality. Adding two collision terms also requires a
nonoverlapping process decomposition; it can otherwise double count encounters.

\section{Research questions and evidence}

\textbf{Geometry and nullspace.} Begin with the proven finite-quadrature spectrum
\eqref{eq:uniformspectrum}. Track its changes under a valid mirror, toroidal or
dipole field. Separate collision null modes from decay generated through coupling
to Hamiltonian transport. A small positive eigenvalue on one finite grid is not a
continuum gap. Vary all resolutions, domains and velocity tails.

\textbf{Localization.} Compare the local kernel with a specified symmetric
finite-range kernel, using a derived normalization as the range shrinks.
Boundary renormalization must retain pair symmetry. A Gaussian-range surrogate
remains a surrogate unless independent scattering physics supports it.

\textbf{Constrained dipole states.} Ref.~\cite{S23} studies maximum-entropy dipole
states and additional adiabatic constraints. Here the result sought is a dynamical
relaxation law, an extra invariant, a loss of mixing, or an independently measured
lifetime. Reproducing \eqref{eq:density} is verification rather than discovery.

For any geometry-amplitude scaling, compare competing fits and inspect projector
degeneracies before using regular perturbation theory. For steady-state parameter
derivatives specify which particle, energy and marginal constraints are held
fixed; otherwise a nullspace can make the derivative underdetermined.

\subsection*{Executed starter evidence}

The supplied NumPy example uses 216 velocity nodes. Its normalized exact decay
rate is $0.11742603550295856$. During package preparation, the direct pair
contraction agreed with \eqref{eq:uniformreduction} to $3.63\times10^{-16}$ in
relative norm. Homogeneous and density-like collision actions were zero in that
run. The entropy-weighted nonzero spectrum agreed with $-\lambda$ to
$3.20\times10^{-16}$ in absolute error. These are algebra checks, not time-series
PDE benchmarks. The independent test file passed five tests, including tilted
fields, both charge signs and parallel wave vectors.

The script, inputs and JSON are in the starter repository. Roundoff and elapsed
time vary by machine. The earlier supplied script and its output are preserved
outside the public starter for provenance. No nonuniform nonlinear simulation,
physical coefficient calibration or new dipole result is claimed here.

The implementation handoff and validation matrix prescribe the remaining tests:
conservation, positive entropy production, positivity, limiting cases,
independent comparisons, numerical convergence and matched-error performance.
Update this section with measured results as those stages are completed.

\bibliographystyle{unsrt}
\bibliography{references}
\end{document}
